(1) 对 \(b\) 模偏迹，交叉 \(n\ne m\) 因正交性消失，得
\[
\rho_a=(1-\lambda^2)\sum_{n=0}^\infty\lambda^{2n}|n\rangle\langle n|.
\]
这是几何热态，\(\bar n=\lambda^2/(1-\lambda^2)=\sinh^2r\)。熵从 \(-\sum_np_n\ln p_n\) 的几何级数求得
\[
\boxed{S=(\bar n+1)\ln(\bar n+1)-\bar n\ln\bar n}.
\]
总态是纯态，两模约化熵相等，所以该熵就是两模的纠缠熵，而非来自外界耗散的热熵。

(2) 取双模压缩变换 \(a_{\rm out}=\cosh r\,a_{\rm in}+\sinh r\,b_{\rm in}^\dagger\)，\(b_{\rm out}=\cosh r\,b_{\rm in}+\sinh r\,a_{\rm in}^\dagger\)，输入双真空。于是
\(x_a-x_b=e^{-r}(x_{a0}-x_{b0})\)，
\(p_a+p_b=e^{-r}(p_{a0}+p_{b0})\)。两个独立真空的差、和方差均为1，因此
\[
\boxed{\mathrm{Var}(x_a-x_b)=\mathrm{Var}(p_a+p_b)=e^{-2r}}.
\]
Duan和式 \(2e^{-2r}<2\)（对 \(r>0\)）违反可分态下界2，构成连续变量双模纠缠见证；并不代表任何一个模式本身的两个共轭正交量同时被压缩。

(3) Schmidt振幅 \(\sqrt{p_n}=\sqrt{1-\lambda^2}\lambda^n\)，纯双体态部分转置的迹范数等于
\[
\|\rho^{T_b}\|_1=(\sum_n\sqrt{p_n})^2
=\frac{1-\lambda^2}{(1-\lambda)^2}
=\frac{1+\lambda}{1-\lambda}=e^{2r}.
\]
自然对数定义的对数负性因此为 \(\boxed{E_{\mathcal N}=\ln\|\rho^{T_b}\|_1=2r}\)。它和第一问的熵都随 \(r\) 增长，但数值、可加性和对混合态的含义不同，不能将二者直接等同。

(4) 真空损耗通道可写 \(a_{\rm det}=\sqrt\eta\,a+\sqrt{1-\eta}\,v_a\)，对 \(b\) 类似，其中 \(v_a,v_b\) 是独立真空。EPR组合的方差是信号部分与真空噪声之和：
\[
\boxed{V_{\rm EPR,out}=(1-\eta)+\eta e^{-2r}},\qquad
\bar n_{\rm out}=\eta\sinh^2r.
\]
Duan和式为 \(2[(1-\eta)+\eta e^{-2r}]<2\) 对任何 \(\eta>0,r>0\) 都仍严格成立，所以纯真空无热噪损耗可减弱却不会在此对称理想模型中产生有限效率的彻底不可分性消失；当 \(\eta\to0\) 时和式趋向2、纠缠检测余量趋零。
